% !TEX root = ../../../main.tex
% Appendix entry: interrogation.trains -- "Trains" (Appendix A).
\entryfield{Algorithm}
The adapter (\code{_run} in \code{Adapters/interrogation_trains.py}) maps the chain onto
\code{find_trains} in \code{Working/interrogation/trains.py}. It never reads the signal: it
works on event times only.

\emph{Event times.} With \emph{anchor} = \texttt{onset} and an upstream \code{onset_idx}
column (written by \blockref{interrogation.event_shape}), an event's time is its onset;
otherwise it is its span start. Either is divided by $f_s$ to give seconds. An event whose
onset is NaN gets no time, is in no train, and is counted in \code{n_unanchored}.

\emph{Groups.} If the input has a \code{group} column, events are split by it and trains are
found separately in each group (\code{Working/interrogation/intervals.py}: an interval is only
meaningful inside one group). Train ids are renumbered so they are unique across groups. With
no \code{group} column the whole span is one group.

\emph{Per group}, \code{find_trains} sorts the times $t_1 \le \dots \le t_n$ and forms the
gaps $g_i = t_{i+1} - t_i$. A train is always a run of consecutive events in time order; the
three rules differ only in which gaps join a run.

\begin{description}
\item[\texttt{logisi}] (\code{logisi_trains}, \code{logisi_cut}; after
  \citet{pasquale2010selfadapting} as restated by \citet{cotterill2016comparison}, since the
  2010 full text was not reachable).
  \begin{enumerate}
  \item Fewer than 10 gaps: refuse, label every event $-1$ and say so in \code{refused}.
  \item Take $\ell_i = \log_{10} g_i$ over the positive gaps only (zero gaps, from
    coincident events, have no logarithm; they are left out of the density but still count
    as below any cut). Fit a Gaussian kernel density with Silverman's robust bandwidth
    \[ h = 0.9\,\min\!\bigl(s, \mathrm{IQR}/1.349\bigr)\, m^{-1/5}, \]
    $s$ the sample SD (ddof 1) of the $m$ log-gaps, using $s$ alone when the IQR is 0.
    Evaluate it on 512 points from $\min\ell - 3h$ to $\max\ell + 3h$. If all gaps are
    equal there is no density and no trains.
  \item Keep only peaks whose prominence is at least 5\,\% of the density's maximum
    (\code{PEAK_PROMINENCE_FRAC}).
  \item Candidate intra-train peaks: with \emph{max\_core\_gap\_s} = 0 (auto), every
    prominent peak, left to right; with a ceiling $c > 0$, only the largest peak at or
    below $\log_{10} c$ (none: no trains).
  \item For each candidate $p_1$ and each prominent peak $p_2 > p_1$, take the minimum $m$ of
    the density between them and the void
    \[ v = 1 - \frac{g(m)}{\sqrt{g(p_1)\,g(p_2)}} . \]
    The first pair with $v > 0.7$ (\code{VOID_THRESHOLD}, fixed, not a setting) sets the
    cut $\tau = 10^{\text{grid}[m]}$ seconds. No such pair: no trains, with the deepest void
    and the reason in words.
  \item Auto mode, or a cut not above the ceiling: a train is a run of at least
    \emph{min\_events} events joined by gaps $g_i < \tau$ (strictly).
  \item A ceiling $c$ with $\tau > c$: \code{core_and_extend}. Cores are runs of at least
    \emph{min\_events} events with gaps $< c$; a core grows over any neighbour whose gap to
    it is $< \tau$, repeated until nothing changes; members joined by gaps $< \tau$ form one
    train.
  \end{enumerate}
\item[\texttt{maxinterval}] (\code{maxinterval_trains}; the benchmark's first choice in
  \citet{cotterill2016comparison}). Walking the sorted events: a train begins at the first
  gap $\le$ \emph{mi\_begin\_gap\_s} and continues while gaps are $\le$
  \emph{mi\_end\_gap\_s}. Consecutive trains whose inter-train interval (first event of the
  next minus last event of this one) is $<$ \emph{mi\_min\_ibi\_s} are merged, taking the
  events between them. Then a train lasting less than \emph{mi\_min\_duration\_s} (last minus
  first event) or holding fewer than \emph{min\_events} events is removed. The defaults
  (60, 120, 180, 0\,s) are placeholders for 1\,Hz fungal data; the neuronal values were
  0.17, 0.3, 0.2 and 0.01\,s.
\item[\texttt{adamatzky}] (\code{adamatzky_words}; \citet{adamatzky2022language}). With
  $\bar g$ the mean gap of the group and $\theta = \text{\emph{theta\_mult}}\cdot\bar g$,
  consecutive events share a word while $g_i \le \theta$ (not strict). Every event is in a
  word, words of one event included; a word of at least \emph{min\_events} events is a
  train. The paper's $a$ is a species' mean interspike interval; here it is the mean gap of
  one group in one span.
\end{description}

\emph{Summaries} (all rules). Trains are numbered $0,1,\dots$ in time order;
\code{train_size} and \code{position_in_train} (0-based) follow. \code{lz_complexity} is the
Kaspar--Schuster LZ76 phrase count of the sequence of train sizes (one symbol per train; for
\texttt{adamatzky}, of every word's length), normalised by $n/\log_2 n$ with the routine
\code{_lz76_complexity} from \code{Working/Detection/analysis/dehshibi_detection_analysis.py};
\code{None} below two trains. With more than two symbols it is not bounded by 1.
\code{artefact_warning} warns when the median within-train gap lies within 25\,\% of 200\,s,
the $\sim$0.005\,Hz cycle seen on no-fungus control channels \citep{buffi2025electrical}.

\code{gap_before_s} is not computed here: it is \code{measure} from
\code{Adapters/interrogation_intervals.py}, imported, so it equals the
\code{interval_before_s} that \blockref{interrogation.intervals} would report for the same
anchor and groups.

With \emph{write\_group} on, \code{write_group_columns}
(\code{Working/interrogation/intervals.py}) sets \code{group} to the train id (or
\code{<upstream>/<train id>}, keeping the upstream key in \code{group_upstream}). Every event
in no train shares the key \texttt{-1}.

\entryfield{Assumptions}
\begin{itemize}
\item \emph{Units.} Only times: samples divided by $f_s$, so seconds. The signal's unit
  never enters. The gap settings are in seconds and so mean the same at any sampling rate.
\item \emph{Missing values.} The block does not see the trace, so a NaN stretch in the
  recording is invisible to it: no events are found there, and the gap across it is just a
  long gap. That gap ends a train under every rule, and under \texttt{adamatzky} it also
  raises the mean gap and so $\theta$. Events with a NaN onset are left out (above).
\item \emph{Edges of the span.} No edge handling: the first event has no gap before it
  (\code{gap_before_s} NaN) and the last none after. A train cut by the span's end is
  reported with the events the span holds; nothing marks it as truncated.
\item \emph{Event count.} \texttt{logisi} needs at least 11 events in a group. The others run
  on any number: one event gives one word of one and no train; zero events give an empty
  result with ``no events''.
\item \emph{Signal shape.} \texttt{logisi} assumes the log-gap distribution is a mixture of
  two humps. The module docstring records the resolution limit, measured on synthetic data:
  with exponential (irregular) gaps inside trains, the between-train gaps must be about 30
  to 100 times the within-train ones before a void above 0.7 can exist.
\end{itemize}

\entryfield{Outputs}
A \texttt{SpanSet} with the input's spans, labels, scores and feature columns unchanged, plus
\code{train_id} ($-1$ = none), \code{train_size} (0 when none), \code{position_in_train}
($-1$ when none) and \code{gap_before_s}; \code{word_id} under \texttt{adamatzky}; \code{group}
(and \code{group_upstream}) when \emph{write\_group} is on. A column of the same name already
upstream is replaced.

\code{meta} holds: \code{rule}, \code{n_events}, \code{n_gaps}, \code{n_groups},
\code{n_unanchored}, \code{anchor_used}, \code{n_trains}, \code{n_in_trains},
\code{fraction_in_trains}, \code{train_size_distribution}, \code{refused},
\code{no_trains_reason} (only when no train was found), \code{warnings}, \code{rules} (the
printed rule behind every number, \code{rules} in \code{trains.py}) and \code{groups} (each
group's own summary). When there is one group, its \code{cut_s}, \code{void},
\code{intra_peak_s}, \code{density} (the log-gap density, for drawing), \code{lz_complexity},
\code{median_within_gap_s}, and the rule's own fields (\code{theta_s}, \code{mean_gap_s},
\code{word_length_distribution}, \code{n_words}; \code{thresholds_s}) are copied to the top
level; with several, those are \code{None} at the top and a note points to \code{groups}.

\entryfield{Cost}
$O(n \log n)$ in the number of events (a sort per group) plus a 512-point density, and the
iterative \code{core_and_extend} loop when a ceiling is set. The span length enters only
through the number of events. \code{tests/test_trains.py} checks that 20{,}000 events take
well under a second. The estimate (\code{_estimate}) scales with span length as a stand-in.
Not meant for the cluster.

\entryfield{Failure modes}
\begin{itemize}
\item \emph{``No trains'' that means ``not resolvable''} (above). The reason string says
  which: refused, one hump, a second hump below the prominence floor, or valleys too shallow.
\item \emph{A cut at the wrong valley.} In auto mode the leftmost prominent hump is taken as
  the within-train one. A ripple of short gaps (split detections, duplicate onsets) can form
  that hump; the registry advises setting \emph{max\_core\_gap\_s} when the valley lands in
  the wrong place.
\item \emph{Placeholder thresholds.} \texttt{maxinterval}'s four second-valued defaults are
  not fitted to any fungal recording; at 60\,s and 120\,s they will call trains on any
  recording whose events are that dense.
\item \emph{\texttt{adamatzky} always joins something.} With $\theta$ equal to the mean gap,
  about half of all gaps or more are at or below it on most distributions, so this rule
  finds words in Poisson (structureless) timing too. No null is applied.
\item \emph{Unanchored events vanish from the timing.} If Event shape found no fall in some
  events, their onsets are NaN; they are left out, and the gap around them joins their
  neighbours.
\item \emph{Gaps in the recording} split trains silently (above).
\item \emph{A 200\,s rhythm} may be the control-channel artefact; the block warns.
\end{itemize}

\entryfield{Validation}
\code{tests/test_trains.py} checks answers, not only shapes, on synthetic event times:
\texttt{logisi} finds no trains in 25 Poisson series; recovers 20 clear trains (gaps
20--40\,s inside, 2000--4000\,s between) with exactly the right members in 10 seeds, and
again when fewer than half the gaps are within trains; recalls at least 93\,\% (300 events)
and 75\,\% (60 events) of train events when the timing inside is exponential at a
between-to-within ratio of 100; and calls trains in at most 5\,\% of trainless series
(Poisson and 100-fold rate ramps, 30 and 60 events, 100 seeds each). \texttt{maxinterval}
and \texttt{adamatzky} are checked against hand-worked examples, including a gap equal to
$\theta$ joining a word.

On real recordings the evidence is negative: in
\code{docs/prompts/blocks/reports/W1-evaluation.md} the trains template found 0 events in
trains on three recordings at 4--6\,h and three at 24\,h, and 0 to 0.15 on average on their
phase-randomised surrogates, so it was not shipped as a template and the block is marked
experimental. No test compares its trains with hand-marked trains on a fungal recording.

\entryfield{References}
\citet{pasquale2010selfadapting}, as restated by \citet{cotterill2016comparison};
\citet{cotterill2016comparison} for MaxInterval; \citet{adamatzky2022language} for words;
\citet{buffi2025electrical} for the 200\,s control-channel cycle. The adaptations (a kernel
density instead of a binned histogram, the optional ceiling, no fallback threshold, the
prominence floor, the 10-gap minimum) are each argued with measurements in the docstring of
\code{Working/interrogation/trains.py}; the sources are summarised in
\code{docs/research/analysis-blocks-primary-sources.md}, section 8.

\entryfield{Tests}
\code{tests/test_trains.py} (the core, with the answer checks above);
\code{tests/test_adapter_trains.py} (the four columns, upstream onset as anchor, trains never
crossing a group, \emph{write\_group} with and without an upstream group, the refusal below
ten gaps, an empty span set, Poisson spans giving no trains, JSON-safe meta, linear time over
many groups); \code{tests/test_template_event_trains.py} (the template chain finds four
planted trains end to end).
